% GENERATED FILE. Edit canonical lessons, then run npm run generate:books.
% canonical-source-hash: fnv1a64:200d42589fd4260c
% canonical-lessons: GE-C295-panne, GE-C295-pannendienst, GE-C295-abschleppen, GE-C295-beschaedigen, GE-C295-zerstoeren

\chapter{Eine Panne}
\label{ch:ge-eine-panne}
\hlchaptermodality{295}

\begin{chapteropening}
\textbf{By the end of this chapter:} \emph{I can report a breakdown and damage: a breakdown, the breakdown service, to tow away, to damage and to destroy.}

Complete the last lesson of chapter 295: I can report a breakdown and damage: a breakdown, the breakdown service, to tow away, to damage and to destroy.
\end{chapteropening}

\section[die Panne]{die Panne — the breakdown; a mishap}

\label{lesson:GE-C295-panne}



\begin{quote}
Before the new one: say the German for the high water, then the German for to collapse.
\end{quote}

\subsection*{You'll want to know: die Panne}
\textbf{die Panne} — \textquotedblleft{}the breakdown; a mishap\textquotedblright{}.

\begin{quote}
Auf der Autobahn hatten wir eine Panne.
\end{quote}

\begin{quote}
\emph{We had a breakdown on the motorway.}
\end{quote}

Plural: \emph{die Pannen}. Also a small mishap: \emph{eine kleine Panne}.

\begin{cousinweb}
\emph{Panne} was borrowed from French \emph{panne} \textquotedblleft{}breakdown\textquotedblright{} (EtymDB).
\end{cousinweb}

\begin{grammarlens}[title={What you've built}]
One more word for telling, explaining and describing at B1.
\end{grammarlens}

\subsection*{Your turn}
\begin{itemize}
\raggedright
  \item \emph{Say it:} \emph{die Panne}
  \item \emph{Say it:} \emph{die Panne}, once more
  \item \emph{Say it:} a sentence of your own with \emph{die Panne}
  \item \emph{Recall:} say the German for the high water, then the German for to collapse, then say \emph{die Panne} again
\end{itemize}

\begin{tcolorbox}[breakable,colback=teal!4,colframe=teal!35!black,title={Before you move on}]
What does die Panne mean? (\textquotedblleft{}The breakdown; a mishap\textquotedblright{}.) Say it once more.
\end{tcolorbox}
\section[der Pannendienst]{der Pannendienst — the breakdown service}

\label{lesson:GE-C295-pannendienst}



\begin{quote}
Before the new one: say the German for to collapse, then the German for the breakdown.
\end{quote}

\subsection*{You'll want to know: der Pannendienst}
\textbf{der Pannendienst} — \textquotedblleft{}the breakdown service\textquotedblright{}.

\begin{quote}
Der Pannendienst kam nach einer Stunde.
\end{quote}

\begin{quote}
\emph{The breakdown service came after an hour.}
\end{quote}

Plural: \emph{die Pannendienste}. \emph{Der Dienst} — service.

\begin{cousinweb}
\emph{Pannendienst} joins \emph{Panne}, borrowed from French \emph{panne} \textquotedblleft{}breakdown\textquotedblright{} (EtymDB), and the inherited \emph{Dienst} \textquotedblleft{}service\textquotedblright{}.
\end{cousinweb}

\begin{grammarlens}[title={What you've built}]
One more word for telling, explaining and describing at B1.
\end{grammarlens}

\subsection*{Your turn}
\begin{itemize}
\raggedright
  \item \emph{Say it:} \emph{der Pannendienst}
  \item \emph{Say it:} \emph{der Pannendienst}, once more
  \item \emph{Say it:} a sentence of your own with \emph{der Pannendienst}
  \item \emph{Recall:} say the German for to collapse, then the German for the breakdown, then say \emph{der Pannendienst} again
\end{itemize}

\begin{tcolorbox}[breakable,colback=teal!4,colframe=teal!35!black,title={Before you move on}]
What does der Pannendienst mean? (\textquotedblleft{}The breakdown service\textquotedblright{}.) Say it once more.
\end{tcolorbox}
\section[abschleppen]{abschleppen — to tow away}

\label{lesson:GE-C295-abschleppen}



\begin{quote}
Before the new one: say the German for the breakdown, then the German for the breakdown service.
\end{quote}

\subsection*{You'll want to know: abschleppen}
\textbf{abschleppen} — \textquotedblleft{}to tow away\textquotedblright{}.

\begin{quote}
Das Auto wurde abgeschleppt.
\end{quote}

\begin{quote}
\emph{The car was towed away.}
\end{quote}

Separable: \emph{Sie schleppen das Auto ab}.

\begin{cousinweb}
\emph{abschleppen} is built on \emph{schleppen} \textquotedblleft{}to drag\textquotedblright{}, Middle High German \emph{slepen}, which came from Middle Low German \emph{slēpen} \textquotedblleft{}to drag\textquotedblright{} (EtymDB).
\end{cousinweb}

\begin{grammarlens}[title={What you've built}]
One more word for telling, explaining and describing at B1.
\end{grammarlens}

\subsection*{Your turn}
\begin{itemize}
\raggedright
  \item \emph{Say it:} \emph{abschleppen}
  \item \emph{Say it:} \emph{abschleppen}, once more
  \item \emph{Say it:} a sentence of your own with \emph{abschleppen}
  \item \emph{Recall:} say the German for the breakdown, then the German for the breakdown service, then say \emph{abschleppen} again
\end{itemize}

\begin{tcolorbox}[breakable,colback=teal!4,colframe=teal!35!black,title={Before you move on}]
What does abschleppen mean? (\textquotedblleft{}To tow away\textquotedblright{}.) Say it once more.
\end{tcolorbox}
\section[beschädigen]{beschädigen — to damage}

\label{lesson:GE-C295-beschaedigen}



\begin{quote}
Before the new one: say the German for the breakdown service, then the German for to tow away.
\end{quote}

\subsection*{You'll want to know: beschädigen}
\textbf{beschädigen} — \textquotedblleft{}to damage\textquotedblright{}.

\begin{quote}
Der Ball beschädigte das Fenster.
\end{quote}

\begin{quote}
\emph{The ball damaged the window.}
\end{quote}

Regular: \emph{hat beschädigt}. Built on \emph{der Schaden}.

\begin{cousinweb}
\emph{beschädigen} is built on \emph{Schaden} \textquotedblleft{}damage\textquotedblright{}, which is usually traced to Old High German \emph{scado} \textquotedblleft{}damage\textquotedblright{}, from Proto-Germanic \emph{skaþô} (EtymDB).
\end{cousinweb}

\begin{grammarlens}[title={What you've built}]
One more word for telling, explaining and describing at B1.
\end{grammarlens}

\subsection*{Your turn}
\begin{itemize}
\raggedright
  \item \emph{Say it:} \emph{beschädigen}
  \item \emph{Say it:} \emph{beschädigen}, once more
  \item \emph{Say it:} a sentence of your own with \emph{beschädigen}
  \item \emph{Recall:} say the German for the breakdown service, then the German for to tow away, then say \emph{beschädigen} again
\end{itemize}

\begin{tcolorbox}[breakable,colback=teal!4,colframe=teal!35!black,title={Before you move on}]
What does beschädigen mean? (\textquotedblleft{}To damage\textquotedblright{}.) Say it once more.
\end{tcolorbox}
\section[zerstören]{zerstören — to destroy}

\label{lesson:GE-C295-zerstoeren}



\begin{quote}
Before the new one: say the German for to tow away, then the German for to damage.
\end{quote}

\subsection*{You'll want to know: zerstören}
\textbf{zerstören} — \textquotedblleft{}to destroy\textquotedblright{}.

\begin{quote}
Das Feuer zerstörte das ganze Haus.
\end{quote}

\begin{quote}
\emph{The fire destroyed the whole house.}
\end{quote}

\emph{Die Zerstörung} — the destruction.

\begin{cousinweb}
\emph{zerstören} is built on \emph{stören} \textquotedblleft{}to disturb\textquotedblright{}, Middle High German \emph{stören} (EtymDB), with \emph{zer-} \textquotedblleft{}apart, to pieces\textquotedblright{}.
\end{cousinweb}

\begin{grammarlens}[title={What you've built}]
One more word for telling, explaining and describing at B1.
\end{grammarlens}

\subsection*{Your turn}
\begin{itemize}
\raggedright
  \item \emph{Say it:} \emph{zerstören}
  \item \emph{Say it:} \emph{zerstören}, once more
  \item \emph{Say it:} a sentence of your own with \emph{zerstören}
  \item \emph{Recall:} say the German for to tow away, then the German for to damage, then say \emph{zerstören} again
\end{itemize}

\begin{tcolorbox}[breakable,colback=teal!4,colframe=teal!35!black,title={Before you move on}]
What does zerstören mean? (\textquotedblleft{}To destroy\textquotedblright{}.) Say it once more.
\end{tcolorbox}
